\chapter{面向少目标域样本的域适应SAR目标识别方法}
% \chapter{面向少实测样本的域适应SAR目标识别方法}

\section{引言}
已有的研究在无监督域适应SAR目标识别方面已经取得了可观的成果，
然而现有方法仍依赖大量无标签的目标域数据参与模型训练。比如针对像素级的域对齐问题现有研究主要依赖CycleGAN进行像素到像素的转换，
其包含两个生成器分别负责源域-目标域、目标域-源域转换，以及分别用于判别目标域、源域样本真伪的两个判别器，同时通过循环一致性损失约束转换的可逆性与合理性；
而判别器的训练需要足量目标域真实数据来学习目标域数据的完整统计分布，只有获取大量目标域数据，
判别器才能充分掌握目标域的真实数据特征，从而精准区分生成的伪样本与真实样本，为生成器提供有效的优化方向，
引导生成器生成贴合目标域特性的图像；若目标域数据量不足，判别器无法学习到目标域的完整分布，会出现判别能力退化的问题，
既无法为生成器提供有效的优化依据，也会导致循环一致性损失难以约束出符合目标域真实特性的跨域转换结果，
最终无法实现理想的跨域图像对齐效果。

而针对特征层面的域对齐现有方法仍然主要是基于域对抗训练比如DANN。
其由特征提取器、标签分类器与域判别器构成，
其核心目标是通过域判别器与特征提取器的对抗训练，迫使特征提取器对齐源域和目标域的特征，
其中域判别器的核心任务是判断输入特征来自源域还是目标域，这一过程需要足量的目标域数据（实测SAR图像）
来学习目标域数据的完整统计分布，包括实测SAR图像特有的回波散射特性、斑点噪声分布、背景杂波的统计规律等；
若目标域数据量不足，域判别器无法充分掌握目标域的完整数据分布，会出现判别能力退化的问题，
既无法精准区分源域与目标域的特征，也无法为特征提取器提供有效的优化梯度信号，
无法引导特征提取器学习到适配目标域的域不变特征；同时，DANN的域对齐是面向全局特征分布的对齐，
需要足量目标域样本覆盖目标域的分布多样性，
少量数据会导致样本分布存在偏差，无法支撑实现充分的域对齐，最终导致模型在目标域的识别性能下降。

针对上述问题，本文提出了一种基于像素级域对齐和类混淆损失（Class Confusion Loss，CC-Loss）的少目标域样本无监督域适应方法，
使用一种基于对比学习的非配对式图像翻译网络替换常用的CycleGAN以实现少目标域样本条件下的像素级对齐，
并舍弃了基于对抗训练实现特征对齐的范式，转而基于类混淆损失优化目标域自身的判别边界，
旨在提升现有无监督域适应方法在目标域数据受限条件下的适用性与鲁棒性，
减少对目标域数据的依赖，从而提升模型在目标域的识别性能。

% \section{基于像素级域对齐和类混淆损失的少目标域样本域适应方法}
\section{基于像素级域对齐和类混淆损失的域适应SAR目标识别方法}

\subsection{方法概述}

\begin{figure}[htp]
\centering
\includegraphics[width=\textwidth, keepaspectratio]{域自适应框图2}
\caption{基于像素级域对齐和类混淆损失的少目标域样本域适应方法框图}
\end{figure}

图2.1展示了所提出方法的整体框架。整体分为两个模块：像素级域对齐和类混淆域适应模块。
首先，在像素级域对齐模块中使用了一种新颖的适用于少目标域样本条件的非配对式图像到图像翻译模型。
这一模块旨在从SAR图像的像素层面缩小源域和目标域之间的分布差距，其关注图像内部的图像块而不是整幅图像，
可以实现更精确的目标特性保留，并在有限数据的情况下进行高效学习。
随后，将转换后的中间域图像作为新的源域数据输入类混淆域适应模块完成域适应目标识别模型的训练，
基于类混淆的域适应图像分类无需显式对齐源-目标域特征分布，
核心是通过优化目标域自身的类别混淆实现适配：先利用带标注的源域数据，
以交叉熵损失（Cross-Entropy Loss）训练 “特征提取器+分类器” 基础模型；再针对无标注的目标域图像，经特征提取器得到特征后，
通过概率校准缓解分类器过度自信，结合不确定性重加权生成类别混淆矩阵。最后以类别混淆损失（Class Confusion Loss，CC-Loss）为核心优化目标：
一方面最小化目标域跨类别混淆以清晰化类别边界，另一方面强制原始图像与增强图像的混淆矩阵一致以提升鲁棒性，
最终让特征提取器学到目标域专属、类别可区分且抗扰动的通用特征，省去了复杂的特征空间对齐过程，减少了对目标域数据的依赖。



\subsection{像素级域对齐模块}

先前的研究中发现，图像到图像翻译模型可以显著缩小像素级的域差距\cite{chen2022pixel,youk2023transformer}。
在无监督域适应任务中，由于不同域之间的图像对应关系未知，只能使用非配对式的图像翻译模型。
而当目标域可用于训练参考的数据量不足时，必须选择能充分利用信息、在小样本条件下有效工作的图像翻译模型。
针对此，本研究将对比学习和生成对抗模型进行结合，引入了一种新颖的非配对式图像翻译模型（Contrastive Unpaired Translation，CUT）作为像素级域对齐模块。

\begin{figure}[htp]
\centering
\includegraphics[width=0.9\textwidth, keepaspectratio]{cut}
\caption{基于对比学习的像素级域对齐模块CUT结构示意图}
\end{figure}


如图2.2所示，输入两组来自不同域且不配对的SAR图像：$X = \{ x \in \mathcal{X} \}$，$Y = \{ y \in \mathcal{Y} \}$。CUT的生成器$G$分为编码器和解码器两个部分，CUT首先使用编码器$G_{l,\text{enc}}$将输入图像$X$分解为多个特征图：
\begin{equation}
Z_l = H_l\left(G_{l,\text{enc}}(X)\right)
\end{equation}
其中$l$代表层索引，$H_l$是一个双层的多层感知机网络。接着利用解码器$G_{l,\text{dec}}$将特征图映射为生成图像$\hat{Y} = G(X)$。生成的图像仍要经过编码得到特征编码：
\begin{equation}
\hat{Z}_l = H_l\left(G_{l,\text{enc}}(\hat{Y})\right)
\end{equation}
用于计算块噪声对比估计（PatchNCE）损失。PatchNCE损失用于比较生成图像$\hat{Y}$中的每个图像块
与其在源图像$X$中对应位置图像块，使得对应位置的特征尽可能保持一致，同时将源图像$X$中的其他位置的图像块视为负样本，使得不同位置的特征差距尽可能大。NCE指的是噪声对比估计，这个损失函数鼓励编码器学习一种表示，使得对应图像块在特征空间中更接近。同时，它确保了非对应图像块之间的距离更远，从而最大化互信息。$\text{PatchNCE}$损失具体定义为：
\begin{equation}
L_{\text{PatchNCE}}(G,H,X) = \mathbb{E}_{X \sim \mathcal{X}} \sum_{l=1}^{L} \sum_{s=1}^{S_l} l_{\text{NCE}} \left( \hat{z}_l^s, z_l^s, \hat{z}_l^{S \setminus s} \right)
\end{equation}
其中，$l$和$s$分别代表层索引和层内的空间位置，$\hat{z}_l^s$是查询键所对应的图像块的特征向量，$z_l^{S \setminus s}$是来自负图像块的特征向量的堆叠。而$l_{\text{NCE}}$代表着信息噪声对比估计（$\text{InfoNCE}$）损失函数，其被用于在对比学习中最大化互信息。
对比学习的问题有三个信号组成：查询键、正样本和负样本，要做的就是让查询键和正样本信号相关联和负样本形成对比。
首先将查询键、正样本以及$N$个负样本，分别映射成$K$维向量$\mathbf{v}, \mathbf{v}^+ \in \mathbb{R}^K, \mathbf{v}^- \in \mathbb{R}^{N \times K}$，其中$\mathbf{v}_n^- \in \mathbb{R}^K$表示第$n$个负样本将这些样本归一化至单位球中，防止空间扩张或坍缩。
这样就形成了一个$N+1$的分类问题。$l_{\text{NCE}}$交叉熵损失计算如下，其中$\tau$是比例超参，常被称为温度系数，表示正样本被选中的概率。
\begin{equation}
l_{\text{NCE}}(\mathbf{v}, \mathbf{v}^+, \mathbf{v}^-) = -\log\left[ \frac{\exp(\mathbf{v} \cdot \mathbf{v}^+ / \tau)}{\exp(\mathbf{v} \cdot \mathbf{v}^+ / \tau) + \sum_{n=1}^{N} \exp(\mathbf{v} \cdot \mathbf{v}^- / \tau)} \right]
\end{equation}
在图像对抗生成部分，CUT 仍然是使用 GAN 的判别器对抗损失，本研究中采用是交叉熵损失，保证生成的图像能和目标域的图像尽可能相似，这部分的损失公式如下：
\begin{equation}
L_{\text{GAN}}(G,D,X,Y) = \mathbb{E}_{y \sim Y} \log D(y) + \mathbb{E}_{x \sim X} \log\bigl(1 - D(G(x))\bigr)
\end{equation}
其中，$D$ 是判别器，本文中使用的是两层卷积层结构的二元分类器，$y$ 是来自目标域的真实图像，而 $x$ 是来自源域的输入图像。
最终的优化目标损失函数中还添加了一个一致性损失 $L_{\text{PatchNCE}}(G,H,Y)$ 来最小化目标域图像的变化 $\mathbb{E}_y \| G(y) - y \|_1$，
避免生成器对生成图像目标特性进行不必要的改变。综上所述，总损失包括三部分：
对抗性损失$L_{\text{GAN}}(G,D,X,Y)$、对比损失$L_{\text{PatchNCE}}(G,H,X)$和一致性损失$L_{\text{PatchNCE}}(G,H,Y)$，总的损失函数为三项损失的和，具体公式如下：
\begin{equation}
L_{\text{cut}} = L_{\text{GAN}}(G,D,X,Y) + L_{\text{PatchNCE}}(G,H,X) + L_{\text{PatchNCE}}(G,H,Y)
\end{equation}
CUT关注图像内部的图像块而不是整幅图像，可以实现更精确的目标特性保留，并在有限数据的情况下进行高效学习。
CUT基于对比学习的思想，对每个生成的图像块将其在源图像中对应位置图像块作为正样本，并将源图像中不同位置的图像块作为负样本，
以实现保留图像空间结构特征的同时差异化风格。其正样本和负样本均来自同一图像，减少了大型数据集的需求，并简化了学习过程。
所提出的模型由单个生成器和判别器组成，因此相对其他模型参数量更小，更不容易过拟合，因此整体更适用于少目标域样本条件下的域适应任务。


\subsection{类混淆域适应模块}

在SAR目标识别领域现有的无监督域适应算法主要使用域对抗训练来实现源-目标域的特征对齐，
然而这需要大量的目标域数据用于对抗学习，当目标域数据量不足时，跨域效果会显著下降。
为此，本研究舍弃了特征对齐范式，转而针对少目标域样本条件基于类混淆损失和随机数据增强（RandAugment）设计了类混淆域适应模块：
以 “类混淆” 为核心归纳偏置，通过估计并最小化目标域中分类器的类间混淆程度，同时增强混淆矩阵在数据增强下的一致性，
实现跨场景的域自适应性能提升。类混淆损失仅依赖目标域分类器预测结果，不依赖特定标签集或域配置，具备极强的通用性。

类混淆块（Class Confusion Block, CC-Block）是类混淆域适应模块的核心组件，负责准确量化目标域中的类混淆程度，包含概率校准、
不确定性加权和类别归一化三个关键步骤：

\begin{figure}[htp]
\centering
\includegraphics[width=0.8\textwidth, keepaspectratio]{ccblock}
\caption{类混淆块（CC-Block）结构示意图}
\end{figure}

（1）概率校准：为缓解分类器在目标域中过度自信的问题，CC-Block采用温度缩放技术，对分类器的 logit 输出进行校准，得到更可靠的预测概率：
\begin{equation}
\hat{Y}_{ij} = \frac{\exp\left(Z_{ij} / T\right)}{\sum_{j'=1}^{K} \exp\left(Z_{ij'} / T\right)}
\end{equation}
其中，$Z_{ij}$ 为分类器输出层（softmax前）的logit值，$T$ 为温度缩放超参数（$T=1.0$ 时即为标准softmax），
$K$ 为源域类别数，$\hat{Y}_{ij}$ 表示第 $i$ 个目标样本被预测为第 $j$ 类的校准后概率。

（2）不确定性加权：不同目标样本对类混淆量化的贡献存在差异：预测概率呈多峰分布的样本（分类器在多个类别间犹豫）更能体现类混淆，
而均匀分布或单峰分布的样本贡献较低。为提升类混淆估计的鲁棒性，CC-Block基于预测概率计算每个样本的不确定性，求得样本预测熵$H\left(\hat{y}_{i\cdot}\right)$：
\begin{equation}
H\left(\hat{y}_{i\cdot}\right) = -\sum_{j=1}^{K} \hat{Y}_{ij} \log \hat{Y}_{ij}
\end{equation}
其中，$\hat{y}_{i\cdot}$ 表示第 $i$ 个样本的预测概率分布,并通过 softmax 变换赋予高贡献样本更大权重：
\begin{equation}
W_{ii} = \frac{B\left(1 + \exp\left(-H\left(\hat{y}_{i\cdot}\right)\right)\right)}{\sum_{i'=1}^{B} \left(1 + \exp\left(-H\left(\hat{y}_{i'\cdot}\right)\right)\right)}
\end{equation}
其中，$B$ 为批大小，$W_{ii}$ 为第 $i$ 个样本的权重，权重矩阵 $W$ 为对角矩阵，且满足 $\sum_{i=1}^{B} W_{ii}=B$（保证平均权重为1）。

（3）类别归一化：批处理场景下易存在类别不平衡问题，导致类混淆估计偏差，为消除类别不平衡对类混淆估计的影响，
CC-Block对类别混淆矩阵进行归一化处理，确保每个类别在混淆矩阵中的贡献均等，首先基于加权后的预测概率初步计算类混淆：
\begin{equation}
C_{jj'} = \hat{y}_{\cdot j}^{\top} W \hat{y}_{\cdot j'}
\end{equation}
其中，$\hat{y}_{\cdot j}$ 为所有样本对第 $j$ 类的预测概率向量，$C_{jj'}$ 表示第 $j$ 类与第 $j'$ 类的初步混淆值。
随后对混淆矩阵按行归一化，得到最终的类别归一化混淆矩阵：
\begin{equation}
\hat{C}_{jj'} = \frac{C_{jj'}}{\sum_{j''=1}^{K} C_{jj''}}
\end{equation}
归一化后，$\hat{C}_{jj'}$ 可解释为“从第 $j$ 类误分类到第 $j'$ 类的概率”，有效缓解类别不平衡影响。

CC-Loss由类混淆最小化和类混淆一致性两部分组成，分别实现“降低类间混淆”和“增强模型鲁棒性”的目标。
类混淆最小化核心目标是最小化跨类（$j \neq j'$）的混淆程度，使分类器对目标域样本的预测更明确。
当跨类混淆最小时，类内一致性自然最大化，定义损失项为：
\begin{equation}
\mathcal{L}_{\text{min}} = \frac{1}{K} \sum_{j=1}^{K} \sum_{j' \neq j} \hat{C}_{jj'}
\end{equation}
其中，$K$ 为类别数，该损失项通过最小化所有跨类混淆值的平均值，引导分类器减少类别间的预测模糊性。

数据增强会引入分布扰动，导致混淆矩阵不稳定。类混淆一致性的目标是基于矩阵级一致性正则化，要求原始数据与增强数据的混淆矩阵保持一致：
首先计算增强数据的类混淆：对目标域样本施加随机数据增强得到 $X^{\text{aug}}$，其校准后预测为 $\hat{Y}^{\text{aug}}$，
加权后类混淆为 $C_{jj'}^{\text{aug}}$，归一化后为 $\hat{C}_{jj'}^{\text{aug}}$。
接着计算其一致性损失项（即最小化L2距离）：
\begin{equation}
\mathcal{L}_{\text{consis}} = \frac{1}{K} \sum_{j=1}^{K} \sum_{j' \neq j} \left(\hat{C}_{jj'} - \hat{C}_{jj'}^{\text{aug}}\right)^2
\end{equation}

整合类混淆最小化和一致性正则化，CC-Loss的最终形式为：
\begin{equation}
\mathcal{L}_{\text{CC}}(\hat{Y}_t, \hat{Y}_t^{\text{aug}}) = \frac{1}{K} \sum_{j=1}^{K} \sum_{j' \neq j} \left( \hat{C}_{jj'} + \left(\hat{C}_{jj'} - \hat{C}_{jj'}^{\text{aug}}\right)^2 \right)
\end{equation}
其中，$\hat{Y}_t$ 和 $\hat{Y}_t^{\text{aug}}$ 分别为目标域原始样本和随机增强样本的通过特征提取器和分类器后得到的校准后预测矩阵。


所提少目标域样本域适应方法总体流程如图2.1所示，
首先，CUT将源域图像转换为更加接近目标域的中间域图像作为新的源域数据$X_s$，以缩小源域与目标域之间的差距，使类混淆在两个域之间更加一致，减小跨域难度。
随机数据增强是一种轻量化自动数据增强方法，摒弃复杂搜索流程，仅通过两个超参数即可实现高效增强，适配各类视觉任务且无需预训练适配。
其核心逻辑为从直方图均衡化、旋转、曝光、改变对比度等14种预设图像变换操作中随机选取N个，为选中操作统一分配强度M并顺序执行，
通过线性映射将M转化为各操作的具体参数。
它用于随机增强目标域图像得到$X_t^{\text{aug}}$，然后和原目标域图像$X_t$计算类别混淆一致性。接着，将得到的源域、目标域和增广图像都输入目标识别模型进行训练。
本章使用ResNet50作为特征提取器，全连接层作为分类器。
源域图像经过特征提取器和分类器后得到预测类别概率矩阵$\widehat{Z}_s$，与标签一起计算源域的交叉熵分类损失，使得模型具备基础的目标识别能力：
\begin{equation}
\mathcal{L}_{\text{CE}} = -\mathbb{E}_{(x_s, y_s) \sim (X_s, Y_s)} \sum_{k=1}^{K}  y_s \log P(y_s = k | x_s)
\end{equation}
其中，$(X_s, Y_s)$ 为带标签的源域数据集，$P(y_s = k | x_s)$ 为分类器对源域样本 $x_s$ 预测为类别 $k$ 的概率。
最后，利用目标域数据和随机增广数据的预测概率矩阵$\widehat{Z}_t$和$\widehat{Z}_t^{\text{aug}}$
计算目标域的类混淆损失 $\mathcal{L}_{\text{CC}}$，最终的训练目标为最小化源域的分类损失和目标域的类混淆损失的之和：
\begin{equation}   
\mathcal{L}_{\text{total}} = \mathcal{L}_{\text{CE}} +  \mathcal{L}_{\text{CC}}
\end{equation}




\section{实验设计与结果分析}

\subsection{实验数据集介绍}
为了验证本章所提方法在跨域SAR目标识别任务中的有效性，本节使用公开的SAMPLE数据集进行实验评测。
SAMPLE数据集包含充足的仿真与实测成对SAR图像，是跨域SAR目标识别任务的研究中最广泛使用的数据集之一。
% SAMPLE数据集表格排版代码
\begin{table}[htb]
  \centering
  % 沿用教程中的minipage环境，宽度设为0.9\linewidth（匹配教程示例）
  \begin{minipage}[t]{\linewidth}
  % 标题格式：中括号内为索引短标题，外部为完整标题（贴合教程规范）
  \caption[SAMPLE数据集十类目标型号与数量]{SAMPLE数据集十类目标型号与数量}
  \label{tab:sample-data}
    % 表格列格式：
    % 第1列：居中+黑体+固定宽度（适配换行）；后10列使用教程中提到的Z列（自动换行+居中）
    \begin{tabularx}{\linewidth}{>{\centering\bfseries\arraybackslash}m{1.5cm}*{5}{>{\centering\arraybackslash}p{0.8cm}}*{5}{>{\centering\arraybackslash}X}}
      \toprule[1.5pt]  % 教程指定的1.5pt粗上框线
      目标类别 & 2S1 & BMP2 & BTR70 & T72 & ZSU23 & M1 & M2 & M35 & M548 & M60 \\
      \midrule[1pt]    % 教程指定的1pt粗中框线
    %   型号 & B01 & 9563 & C71 & 812 & D08 & 0AP00N & MV02GX & T839 & C245HAB & 3336 \\
      型号 & B01 & 9563 & C71 & 812 & D08 & \makecell{0AP\\00N} & \makecell{MV0\\2GX} & T839 & \makecell{C245\\HAB} & 3336 \\
      实测目标图像数量 & 174 & 107 & 92 & 108 & 174 & 129 & 128 & 129 & 128 & 176 \\
      仿真目标图像数量 & 174 & 107 & 92 & 108 & 174 & 129 & 128 & 129 & 128 & 176 \\
      \bottomrule[1.5pt] % 教程指定的1.5pt粗下框线
    \end{tabularx}
  \end{minipage}
\end{table}

% 正文区域的图片排版代码
\begin{figure}[htbp]  % 推荐的浮动参数，给LaTeX排版灵活性
  \centering          % 整体居中
  
  % 第一行：4张子图，用\qquad分隔（教程说明等价于2em间距）
  \subfloat[2S1仿真]{\label{2S1仿真} \includegraphics[width=0.21\textwidth]{2s1_synth_A_elevDeg_016_azCenter_010_22_serial_b01.png}} \hspace{0.8em}
  \subfloat[M1仿真]{\label{M1仿真} \includegraphics[width=0.21\textwidth]{m1_synth_A_elevDeg_016_azCenter_010_18_serial_0ap00n.png}} \hspace{0.8em}
  \subfloat[M35仿真]{\label{M35仿真} \includegraphics[width=0.21\textwidth]{m35_synth_A_elevDeg_016_azCenter_010_62_serial_t839.png}} \hspace{0.8em}
  \subfloat[M548仿真]{\label{M548仿真} \includegraphics[width=0.21\textwidth]{m548_synth_A_elevDeg_016_azCenter_010_63_serial_c245hab.png}} \\[5pt]  % 换行并加5pt垂直间距

  % 第二行：4张子图，保持和第一行一致的宽度/间距
  \subfloat[2S1实测]{\label{2S1实测} \includegraphics[width=0.21\textwidth]{2s1_real_A_elevDeg_017_azCenter_010_22_serial_b01.png}} \hspace{0.8em}
  \subfloat[M1实测]{\label{M1实测} \includegraphics[width=0.21\textwidth]{m1_real_A_elevDeg_017_azCenter_012_18_serial_0ap00n.png}} \hspace{0.8em}
  \subfloat[M35实测]{\label{M35实测} \includegraphics[width=0.21\textwidth]{m35_real_A_elevDeg_017_azCenter_010_62_serial_t839.png}} \hspace{0.8em}
  \subfloat[M548实测]{\label{M548实测} \includegraphics[width=0.21\textwidth]{m548_real_A_elevDeg_017_azCenter_010_63_serial_c245hab.png}} \\

  \caption{SAMPLE数据集在$16^{\circ}$俯仰角下的仿真图像与$17^{\circ}$俯仰角下对应的实测图像}  % 整体大标题
  \label{fig:sample_data}                   % 整体标签，用于交叉引用
\end{figure}



SAMPLE数据集是由美国空军研究实验室（Air Force Research Laboratory，AFRL）发布的公开基准数据集\cite{lewis2019sar}。
该数据集包含10类军用车辆目标，共计1345对配对的仿真与实测SAR图像切片。
其中，仿真数据是利用高精度的CAD模型，通过电磁计算软件模拟生成的，并在几何结构上与实测数据进行了精细对齐。
具体来说，对于给定的CAD模型，模拟合成孔径雷达（SAR）数据的生成主要分为两个步骤：
（1）基于电磁计算，预测目标的电磁散射特性并生成雷达回波；
（2）基于SAR成像算法，对雷达回波进行成像处理，生成复图像。
SAMPLE 数据集中采用的 SAR 成像算法为极坐标格式算法（polar format algorithm）
在成像参数方面，SAMPLE数据集采用X波段HH极化方式，分辨率为0.3m。数据覆盖的俯仰角（Depression angle）范围为$14^{\circ}$至$17^{\circ}$，
方位角（Azimuth angle）范围为$10^{\circ}$至$80^{\circ}$。
表\ref{tab:sample-data}给出了公开发布的SAMPLE数据集中目标的具体型号和数量。
图\ref{fig:sample_data}展示了SAMPLE数据集中部分目标的仿真与实测SAR图像样例。从图中可以看出，尽管仿真图像与实测图像在
目标结构上基本一致，但在纹理细节、噪声特性等方面存在显著差异，这正是域偏移问题的典型体现。
前人已经在SAMPLE数据集上开展了大量的域适应研究工作，积累了丰富的实验结果和对比基准，但少目标域样本域适应问题尚未得到充分探讨。
本文选用SAMPLE数据集作为域适应SAR目标识别的评测基准，
旨在验证所提方法在解决“仿真到实测”域适应问题上的有效性和优越性，尤其是在少目标域数据条件下的域适应鲁棒性。



\subsection{实验设置}


\subsubsection{评价指标}

% 在目标识别任务中，准确率（Accuracy, ACC）与宏平均 F1 分数（Macro F1-score）是评估模型性能的两个基础指标。
% 准确率衡量的是所有预测结果中，预测正确的样本所占的整体比例。因其直观且易于解释，准确率常被用作通用的基准评价指标。
% 然而，在处理类别不平衡数据集时，由于预测结果往往由样本量较大的“多数类”主导，单纯依赖准确率可能会对模型性能产生误导性评价。
% 相比之下，宏平均 F1 分数通过综合考量所有类别的精确率（Precision）与召回率（Recall）来评估模型表现。
% 准确率侧重于整体预测的正确性，而宏平均 F1 分数则更关注每个类别内部精确率与召回率的平衡，确保模型能够兼顾所有类别，避免忽视任何一类（特别是少数类）。
% 因此本文选择同时使用准确率和宏平均 F1 分数两种指标进行综合评价。上述指标的具体计算方式如下：
% \begin{gather}
% \text{Accuracy} = \frac{TP + TN}{TP + TN + FP + FN}  \\
% \text{precision} = \frac{TP}{TP + FP} \\
% \text{recall} = \frac{TP}{TP + FN} \\
% \text{F1-score} = 2 \times \frac{\text{precision} \times \text{recall}}{\text{precision} + \text{recall}} \\
% \text{Macro F1} = \frac{1}{n} \sum_{i=1}^{n} \text{F1-score}_i
% \end{gather}
% 其中，$\mathbf{TP}$（True Positives，真阳性）指模型正确预测为阳性的样本数量；$\mathbf{TN}$（True Negatives，真阴性）表示模型正确预测为阴性的样本数量；
% $\mathbf{FP}$（False Positives，假阳性）代表被错误预测为阳性的样本；而$\mathbf{FN}$（False Negatives，假阴性）则指被错误预测为阴性的样本。
% $\mathbf{Macro\ F1}$（宏F1值）表示$n$个类别中每个类别$i$对应的F1值的平均值。
% 通过同时考察准确率和宏平均 F1 分数，本文能够全面评估所提方法在跨域SAR目标识别任务中的性能表现，
% 确保模型在整体预测正确性的同时，也能兼顾各类别的识别能力。


在目标识别任务中，准确率（Accuracy, ACC）是评估模型性能的基础评价指标。
准确率的核心定义为：在模型输出的所有预测结果中，预测类别与样本真实类别完全一致的样本数量占总样本数量的比例。
因其直观且易于解释，常被用作通用的基准评价指标。本章实验选取准确率作为评价指标，具体计算方式如下：
\begin{gather}
\text{Accuracy} = \frac{TP + TN}{TP + TN + FP + FN} = \frac{N_{\text{correct}}}{N_{\text{total}}}
\end{gather}
其中各符号的物理意义结合SAR目标识别场景解释如下：
${TP}$（True Positives，真阳性）表示模型将真实为目标类的SAR样本正确预测为目标类的数量；
${TN}$（True Negatives，真阴性）表示模型将真实为非目标类的SAR样本正确预测为非目标类的数量；
${FP}$（False Positives，假阳性）表示模型将真实为非目标类的SAR样本错误预测为目标类的数量；
${FN}$（False Negatives，假阴性）表示模型将真实为目标类的SAR样本错误预测为非目标类的数量；
$N_{\text{correct}}$ 表示所有预测正确的样本总数（即 $TP + TN$）；
$N_{\text{total}}$ 表示参与评估的样本总数量（即 $TP + TN + FP + FN$）。

准确率的取值范围为 $[0,1]$，取值越接近1，表明模型的整体预测效果越好；取值越接近0，则说明模型的整体预测错误率越高。
通过准确率指标，本章可直观评估所提方法在跨域SAR目标识别任务中的整体预测正确性，为模型性能提供量化评价基准。







\subsubsection{实验数据设置}

SAMPLE数据集的仿真数据与实测数据之间存在显著的域差异。
为了验证本章所提方法在域适应SAR目标识别任务上的性能，
本节实验使用SAMPLE数据集中的带标签的仿真SAR图像作为源域，无标签的实测SAR图像作为目标域。
所有SAR图像均采用幅度（Amplitude）格式，以16位无符号整数存储，图像尺寸为128×128。
并在输入网络前归一化至 $[0, 1]$ 范围。
为了模拟少目标域样本条件下的域适应任务，
本文分别设置了“all”、“1-shot”、“3-shot”、“5-shot”和“10-shot”五种目标域样本数量的实验场景。
其中，“all”表示使用所有1345张无标签的实测图像作为无监督的目标域参与训练，
而“1-shot”、“3-shot”、“5-shot”和“10-shot”分别表示每个类别仅使用1张、3张、5张和10张实测图像参与训练。
具体实验数据设置如表2.2所示。

% \begin{table}[htb]
%   \centering
%   \begin{minipage}[t]{\linewidth}
%   \caption[SAMPLE数据集少目标域样本域适应实验设置]{SAMPLE数据集少目标域样本域适应实验设置}
%   \label{tab:DA-dataset-settings}
%     \begin{tabularx}{\linewidth}{>{\centering\bfseries\arraybackslash}m{2cm}*{7}{>{\centering\arraybackslash}X}}
%       \toprule[1.5pt]
%       % ==========【关键修改 1/2 核心修改处】==========
%       % multirow{2}{*} 表示：跨行2行、宽度自适应、内容居中，完美实现Dataset在两行中间
%       \multirow{2}{*}{\bfseries 类别} & \multicolumn{6}{c}{\bfseries 训练集} &  {\bfseries 测试集} \\
%       \cmidrule(lr){2-7} \cmidrule(lr){8-8}
%       % ==========【关键修改 2/2 必改处】==========
%       & 源域 & 目标域（all） & 目标域(1shot) & 目标域(3shot)& 目标域(5shot) & 目标域(10shot) \\ 目标域（all） % 这一行的第一列 直接留空即可！！
%       \midrule
%       2S1 & 174 & 174 & 1 & 3 & 5 & 10 & 174 \\
%       BMP2 & 107 & 107 & 1 & 3 & 5 & 10 & 107 \\
%       BTR70 & 92 & 92 & 1 & 3 & 5 & 10 & 92 \\
%       M1 & 129 & 129 & 1 & 3 & 5 & 10 & 129 \\ 
%       M2 & 128 & 128 & 1 & 3 & 5 & 10 & 128 \\
%       M35 & 129 & 129 & 1 & 3 & 5 & 10 & 129 \\
%       M548 & 128 & 128 & 1 & 3 & 5 & 10 & 128 \\
%       M60 & 176 & 176 & 1 & 3 & 5 & 10 & 176 \\
%       T72 & 108 & 108 & 1 & 3 & 5 & 10 & 108 \\
%       ZSU23 & 174 & 174 & 1 & 3 & 5 & 10 & 174 \\
%       \midrule
%       Total & 1345 & 1345 & 10 & 30 & 50 & 100 & 1345 \\
%       \bottomrule[1.5pt]
%     \end{tabularx}
%   \end{minipage}
% \end{table}


\begin{table}[htb]
  \centering
  \begin{minipage}[t]{\linewidth}
  \caption[SAMPLE数据集少目标域样本域适应实验数据设置]{SAMPLE数据集少目标域样本域适应实验数据设置}
  \label{tab:DA-dataset-settings}
    \begin{tabularx}{\linewidth}{>{\centering\bfseries\arraybackslash}m{2cm}*{7}{>{\centering\arraybackslash}X}}
      \toprule[1.5pt]
      \multirow{3}{*}{\bfseries 类别} & \multicolumn{6}{c}{\bfseries 训练集} & \bfseries 测试集 \\
      \cmidrule(lr){2-8}
      & \multirow{2}{*}{\bfseries 源域} & \multicolumn{5}{c}{\bfseries 目标域（无监督）}  & {\bfseries 目标域}  \\ % 修正：增加"测试集"标题
      \cmidrule(lr){3-8} 
      & & all & 1-shot & 3-shot & 5-shot & 10-shot & all\\ % 修正：增加"测试集"标题
      \midrule
      2S1 & 174 & 174 & 1 & 3 & 5 & 10 & 174 \\ % 修正：保留7个数据（无多余列）
      BMP2 & 107 & 107 & 1 & 3 & 5 & 10 & 107 \\
      BTR70 & 92 & 92 & 1 & 3 & 5 & 10 & 92 \\3
      M1 & 129 & 129 & 1 & 3 & 5 & 10 & 129 \\ 
      M2 & 128 & 128 & 1 & 3 & 5 & 10 & 128 \\
      M35 & 129 & 129 & 1 & 3 & 5 & 10 & 129 \\
      M548 & 128 & 128 & 1 & 3 & 5 & 10 & 128 \\
      M60 & 176 & 176 & 1 & 3 & 5 & 10 & 176 \\
      T72 & 108 & 108 & 1 & 3 & 5 & 10 & 108 \\
      ZSU23 & 174 & 174 & 1 & 3 & 5 & 10 & 174 \\
      \midrule
      Total & 1345 & 1345 & 10 & 30 & 50 & 100 & 1345 \\
      \bottomrule[1.5pt]
    \end{tabularx}
  \end{minipage}
\end{table}



\subsubsection{训练参数设置}

本章所提出的两个模块的训练是分开进行的，采用不同的训练参数：
第一步,训练像素级域对齐模块,得到与目标域分布更接近的中间域图像，主要训练参数如表2.3所示。
第二步,在得到的中间域图像基础上,训练类混淆域适应模块,得到具有域适应能力的SAR目标识别模型，主要参数如表2.4所示。


\begin{table}[htb]
  \centering
  \begin{minipage}[t]{0.7\linewidth}
  \caption[像素级域对齐模块训练参数]{像素级域对齐模块训练参数}
  \label{tab:cut-model-hyperparams}
    \begin{tabularx}{\linewidth}{>{\centering\bfseries\arraybackslash}m{5cm} >{\centering\arraybackslash}X}
      \toprule[1.5pt]
      参数 &  \\
      \midrule[1pt]
      优化器 & Adam \\
      初始学习率 & 0.0002 \\
      批次大小 & 1 \\
      网络训练轮数 & 400 \\
      \bottomrule[1.5pt]
    \end{tabularx}
  \end{minipage}
\end{table}


\begin{table}[htb]
  \centering
  \begin{minipage}[t]{0.7\linewidth}
  \caption[类混淆域适应模块训练参数]{类混淆域适应模块训练参数}
  \label{tab:cc-loss-model-hyperparams}
    \begin{tabularx}{\linewidth}{>{\centering\bfseries\arraybackslash}m{5cm} >{\centering\arraybackslash}X}
      \toprule[1.5pt]
      参数 &  \\
      \midrule[1pt]
      初始学习率 & 0.005 \\
      批次大小 & 8 \\
      网络训练轮数 & 20 \\
      每轮迭代次数 & 1000 \\
      优化器 & SGD \\
      动量系数 & 0.9 \\
      权重衰减系数 & 1e-3 \\
      学习率调整策略 & LambdaLR学习率衰减 \\
      温度缩放超参数$T$ & 2.5\\
      \bottomrule[1.5pt]
    \end{tabularx}
  \end{minipage}
\end{table}

在像素级域对齐模块训练中，采用Adam优化器，初始学习率设为0.0002，训练批次大小设为1，总训练轮次为400轮。
在类混淆域适应模块训练中，
温度缩放超参数$T$设置为 2.5，用于CC-Block的温度校准。
模型训练采用随机梯度下降（SGD）优化器，动量因子（Momentum）设为 0.9，权重衰减（Weight Decay）设为 0.001，
以防止模型过拟合。训练批次大小（Batch Size）设定为 8。总训练轮次（Epoch）为 20 轮，每轮包含 1000 次迭代。
学习率调整采用 LambdaLR 设置多项式学习率衰减策略，初始学习率设为 0.005，随迭代步数平滑递减。


硬件与软件环境：
硬件环境搭载Z790 主板、Intel Core i7-13700 CPU 与 NVIDIA GeForce RTX 4090 GPU。软件环境基于 Ubuntu 20.04 操作系统，
利用 PyTorch 1.12.0 深度学习框架实现算法模型的构建、训练与推理测试。



\subsection{对比实验}

为了证明本章所提出方法在域适应SAR目标识别任务中的优越性，以及在少目标域样本条件下的鲁棒性，
本小节基于SAMPLE数据集在各类少目标域样本设置下设计了一系列对比实验。
对比方法选用了五种源自计算机视觉领域的经典无监督域适应算法：
DAN\cite{long2015learning}、DANN\cite{ganin2015unsupervised}、CDAN\cite{long2018conditional}、
DSAN\cite{zhu2020deep}以及MCC\cite{jin2020minimum};
以及五种专为SAR-ATR设计的无监督域适应算法:
PFDA\cite{chen2022pixel}、VSFA\cite{zhang2023vsfa}、Sun等人所提方法\cite{sun2023gradual}、
Shi等人所提方法\cite{shi2024unsupervised}以及MACR\cite{luan2025unsupervised}

具体介绍如下，
DAN（Deep Adaptation Network）\cite{long2015learning}：通过多核极大均值差异MK-MMD度量源域与目标域特征分布差异，并在深度神经网络训练中最小化MK-MMD损失以实现域适应。
DANN（Domain-Adversarial Neural Network）\cite{ganin2015unsupervised}：引入域判别器，通过对抗学习迫使特征提取器对于源域和目标域数据学习到不可区分的特征表示，从而实现域适应。
CDAN（Conditional Domain Adversarial Network）\cite{long2018conditional}：在DANN的基础上，结合类别信息与特征表示进行条件对抗训练，以提升域适应效果。
DSAN（Deep Subdomain Adaptation Network）\cite{zhu2020deep}：通过子域对齐策略，细粒度地匹配源域与目标域的子分布，以增强域适应性能。
MCC（Minimum Class Confusion）\cite{jin2020minimum}：通过最小化类混淆损失来提升目标域上的分类性能，强调类别边界的清晰性，以实现域适应。
PFDA\cite{chen2022pixel}：基于像素级风格迁移模型和特征级域对抗学习，对齐源域和目标域像素级和特征级的分布，以提升SAR目标识别模型的域适应能力。
VSFA\cite{zhang2023vsfa}：基于视觉特征重建与散射拓扑特征融合，利用两阶段损失对齐仿真与实测域的类级别分布，以实现仿真至实测SAR图像的域适应。
Sun等人所提方法\cite{sun2023gradual}：提出渐进式域适应框架，通过特征对齐学习与基于图像相似性的伪标签去噪自训练，逐步缩小域间差异，以实现SAR目标识别模型的域适应。
MACR\cite{luan2025unsupervised}：基于输入–特征–输出三级协同对齐策略，融合风格迁移、对抗学习与知识蒸馏机制，全方位减小多层级域差异，以强化SAR目标识别的跨域类别一致性。
Shi等人所提方法\cite{shi2024unsupervised}：基于域级–类级双粒度对齐，利用梯度加权对抗学习、原型对比学习及属性散射中心辅助机制，实现全局与类级的特征对齐，以提升跨域SAR目标识别性能。

% VSFA\cite{zhang2023vsfa}：融合视觉特征重建与散射拓扑特征，通过两阶段域适应损失在类级别对齐仿真与实测域的SAR图像分布，以实现从仿真SAR图像到实测SAR图像的域适应。
% Sun等人所提方法\cite{sun2023gradual}：
% 构建了渐进式域适应框架，首先通过特征对齐学习域不变特征表示，继而在自训练迭代中引入基于图像相似性的伪标签去噪机制，
% 筛选高置信度目标域样本用于模型微调，逐步缩小仿真域与实测域的分布差异，进而实现SAR目标识别模型的域适应。
% MACR\cite{luan2025unsupervised}：提出输入–特征–输出三级协同对齐策略，实现多层级域适应：输入级风格迁移模型减小像素级域差异；
% 特征级采用对抗学习对齐跨域特征分布；输出级融合跨域知识蒸馏与伪标签过滤，强化类别一致性。
% Shi等人所提方法\cite{shi2024unsupervised}：基于域级和类级双粒度对齐的无监督域自适应方法，通过梯度加权对抗学习实现域级全局对齐、
% 原型网络结合对比学习建模跨域原型关系实现类级对齐、属性散射中心辅助伪标签过滤提升可靠性，显著提升了跨仿真-实测域SAR目标识别性能。

表2.5展示了各方法在SAMPLE数据集不同少目标域样本设置下的实验结果。
除了对比方法以外，本节还设计了两项参考指标：
（1）Source Only：仅使用源域带标签的仿真SAR图像训练ResNet50模型，不进行任何域适应操作。
目标域样本不参与训练，仅作为测试集使用，不影响模型性能表现。
（2）Target Only：直接使用目标域带标签的实测SAR图像训练ResNet50模型。
在“1shot”等少目标域样本条件下，代表仅使用每类1张、3张、5张和10张实测图像进行训练。


\begin{table}[htb]
  \centering
  \begin{minipage}[t]{\linewidth}
  \caption[对比实验结果]{对比实验结果}
  \label{tab:chap2_comparison-results}
    % 修复BUG1：修正表格列格式 正确写法
    % \begin{tabularx}{\linewidth}{>{\centering\bfseries\arraybackslash}m{2cm}>{\centering\bfseries\arraybackslash}m{1.7cm}*{4}{>{\centering\arraybackslash}X}}
    \begin{tabularx}{\linewidth}{>{\centering\bfseries\arraybackslash}m{3cm}*{5}{>{\centering\arraybackslash}X}}
      \toprule[1.5pt]
      \multirow{2}{*}{\bfseries 方法} &  \multicolumn{5}{c}{\bfseries 准确率(\%)}  \\
      \cmidrule(lr){2-6}
      & all & 1-shot & 3-shot & 5-shot & 10-shot \\
      \midrule
      Source Only & 46.29 & 46.29  & 46.29 & 46.29  & 46.29 \\
      Target Only & 99.95 & 16.32 & 32.47 & 54.42 &	80.51 \\
      \midrule
      DAN\cite{long2015learning} & 78.35 & 52.19 & 58.70 & 61.04 & 69.48 \\
      DANN\cite{ganin2015unsupervised} & 83.65 & 48.14 & 52.45 & 59.37 & 71.35 \\
      CDAN\cite{long2018conditional} & 86.96 & 53.48 & 68.20 & 71.74 & 76.12 \\
      MCC\cite{jin2020minimum} & 87.28 & 42.35 & 71.48 & 76.15 & 80.35 \\
      DSAN\cite{zhu2020deep} & 88.95 & 58.83 & 63.87 & 68.12 & 74.15 \\
      \midrule
      PFDA\cite{chen2022pixel} & 96.01 & 不收敛 & 38.17 & 41.58 & 54.37 \\
      VSFA\cite{zhang2023vsfa} & 98.18 &  77.58 & 88.61 & 92.35 & 95.61 \\
      Sun等人方法\cite{sun2023gradual} & 96.46 & 55.47 & 68.28 & 72.67 & 81.14 \\
      Shi等人方法\cite{shi2024unsupervised} & 98.49 & 69.25 & 77.14 & 83.36 & 91.78 \\ 
      MACR\cite{luan2025unsupervised} & 99.27  & 不收敛 & 45.58 & 50.14 & 67.38 \\
      \midrule
      本章方法 &  99.15 & 93.92 & 97.04 & 97.45 & 98.21 \\
      \bottomrule[1.5pt]
    \end{tabularx}
  \end{minipage}
\end{table}







\subsection{消融实验}





\subsection{结果可视化分析}











\section{本章小结}